% GENERATED, edit the French unit or the translation table
\mxunit[3]{terminale-c03}{Limits of sequences}\label{ch:2}

\mxsubtitle{S14, S16, S19, S20}

\begin{keypoint}[Objectives]Distinguish convergence from divergence; calculate a limit; use comparison, the squeeze theorem and monotone convergence.\end{keypoint}

\section{Course notes}
\begin{keypoint}[Definitions]
$u_n\to\ell\in\R$ means that every open interval containing $\ell$ contains all the terms from some index onwards. We say that $(u_n)$ converges to $\ell$.
$u_n\to+\infty$ means that for every real number $A$, there exists an index $N$ such that $n\ge N\Rightarrow u_n\ge A$. $u_n\to-\infty$ means that for every real number $A$, there exists $N$ such that $n\ge N\Rightarrow u_n\le A$.
\end{keypoint}
A sequence may have no limit, even an infinite one: $(-1)^n$ oscillates between $-1$ and $1$. Every convergent sequence is bounded, but the converse is false.
\subsection{Operations and comparisons}
For two convergent sequences with limits $a,b$, the sum tends to $a+b$, the product to $ab$ and the quotient to $\frac{a}{b}$ if $b\ne0$ (the terms of the denominator are then non-zero from some index on). The rules involving infinity are applied by checking signs. The forms $+\infty-\infty$, $0\times\infty$, $\frac{\infty}{\infty}$, $\frac{0}{0}$ do not allow a direct conclusion.
\paragraph{Example.} $\frac{3n^2-n+1}{2n^2+4}=\frac{3-\frac{1}{n}+\frac{1}{n^2}}{2+\frac{4}{n^2}}\to\frac{3}{2}$. For $\sqrt{n^2+n}-n$, multiplying by the conjugate gives $\frac{1}{\sqrt{1+\frac{1}{n}}+1}\to\frac{1}{2}$.
\begin{keypoint}[Theorems]
If $u_n\le v_n$ from some index onwards and $u_n\to+\infty$, then $v_n\to+\infty$. If $a_n\le u_n\le b_n$ from some index onwards and $a_n,b_n\to\ell$, then $u_n\to\ell$ (squeeze theorem).
Every increasing sequence that is bounded above converges. Every decreasing sequence that is bounded below converges. This theorem is accepted without proof.
\end{keypoint}
\paragraph{Proof by comparison.} For every $A$, the terms of $u$ exceed $A$ beyond an index $N_1$. The inequality $v_n\ge u_n$ holds beyond an index $N_2$. For $n\ge\max(N_1,N_2)$, $v_n\ge A$.
\paragraph{Proof: increasing and not bounded above.} For every $A$, the lack of an upper bound provides $N$ such that $u_N>A$. By monotonicity, $n\ge N\Rightarrow u_n\ge u_N>A$, hence $u_n\to+\infty$. Being unbounded alone is not enough.
\subsection{Geometric sequences}
\[q^n\longrightarrow\begin{cases}0&\text{if }-1<q<1,\\1&\text{if }q=1,\\+\infty&\text{if }q>1.\end{cases}\]
If $q\le-1$, $(q^n)$ has no limit: alternation for $q=-1$; for $q<-1$, the even and odd subsequences have opposite signs and unbounded moduli.
\paragraph{Bernoulli and the geometric limit.} Fix $a\ge0$. Let us prove $P(n)$: ``$(1+a)^n\ge1+na$''. At index $0$, both sides equal $1$.
Fix $n\ge0$ and assume $P(n)$ is true. Since $1+a>0$,
\[(1+a)^{n+1}\ge(1+na)(1+a)=1+(n+1)a+na^2\ge1+(n+1)a.\]
This is $P(n+1)$ since $na^2\ge0$. By induction, the inequality is true for every $n\ge0$. If $q>1$, take $a=q-1>0$ and compare with $1+na$. If $0<q<1$, apply this result to $\frac{1}{q}$, then take reciprocals. For $-1<q<0$, $|q^n|=|q|^n\to0$ and the squeeze theorem concludes. The case $q=0$ is handled from index $1$ onwards.
\begin{method}[Recurrence sequence]
For $u_{n+1}=f(u_n)$, establish a bound and monotonicity to prove convergence. If $f$ is continuous at the limit $\ell$, passing to the limit gives $\ell=f(\ell)$ (continuity is studied in chapter~\ref{ch:4}). Solving this equation never proves, on its own, that the sequence converges.
\end{method}


\begin{center}\begin{tikzpicture}\begin{axis}[width=12cm,height=4.6cm,xlabel={$n$},ylabel={$u_n$},xmin=0,xmax=15,ymin=1.5,ymax=5.3,axis lines=left,grid=major,grid style={gray!15}]
\addplot[only marks,mxmint,domain=0:15,samples=16]{2+3/(x+1)};
\addplot[mxink,dashed,domain=0:15]{2};\end{axis}\end{tikzpicture}\par\small The points get closer to $2$ without reaching it: $u_n=2+\frac{3}{n+1}$.\end{center}

\mxnewpage
\section{Research activity}\textbf{A limit is not a last term --- 30 min.}\par Consider the sequence $(u_n)$ defined on $\N$ by $u_n=2+\frac3{n+1}$.
\begin{enumerate}\item Calculate $u_0$, $u_9$ and $u_{99}$. Is there an integer $n$ such that $u_n=2$? Justify.
\item Determine the smallest integer $N$ such that, for every integer $n\ge N$, $|u_n-2|<\num{0,01}$.
\item Let $\varepsilon>0$. Give, in terms of $\varepsilon$, an integer $N$ such that $|u_n-2|<\varepsilon$ for every $n\ge N$.
\item Define on $\N$ the sequence $(v_n)$ by $v_n=2+\frac{(-1)^n}{n+1}$. Prove that it converges to $2$ and that it is not monotonic. What does this example show?\end{enumerate}

\mxnewpage
\section{Progressive exercises}

\mxexerciseitem{1}{Practice}Define on $\N$ the sequences $(u_n)$ and $(v_n)$ by $u_n=\frac{4n-1}{n+2}$ and $v_n=n^2-7n$. Determine their limits, justifying the transformations used.

\mxexerciseitem{2}{Practice}Define $u_n=\sqrt{n+1}-\sqrt n$ for every $n\in\N$. Using the conjugate, determine the limit of $(u_n)$.

\mxexerciseitem{3}{Practice}Define $u_n=\frac{\cos n}{n}$ for every integer $n\ge1$. Bound $u_n$ above and below, then deduce the limit of $(u_n)$.

\mxexerciseitem{4}{Practice}The sequence $(u_n)$ is increasing and satisfies $u_n<5$ for every $n\in\N$. Justify that it converges and state what can be asserted about its limit. Is it necessarily equal to $5$? Justify.

\mxexerciseitem{5}{Practice}Define $u_n=3-2(-\num{0,6})^n$ for every $n\in\N$. Determine the limit of $(u_n)$. Is this sequence monotonic? Justify.

\section{Textbook or past-paper work}\textbf{Barbazo, p. 62, exercises 17 and 24.} In S16: exercise 17, to resolve an indeterminate form by dividing by a power of $n$. At home after S20: exercise 24, to write up an application of the squeeze theorem.

\mxnewpage
\section{Homework}\textbf{Homework 2 --- Approximating a root --- 50 min.}\par Consider $u_0=0$ and, for every $n\in\N$, $u_{n+1}=\sqrt{2+u_n}$. Questions 1, 2 and 4 are to be done after S20; question 3 will be revisited after the lesson on continuity in S31.
\begin{enumerate}\item Prove by induction that all the terms are defined and that $0\le u_n\le2$ for every $n\in\N$.
\item Prove by induction that $u_{n+1}\ge u_n$ for every $n\in\N$. Deduce that $(u_n)$ converges.
\item Using the continuity of the function $x\mapsto\sqrt{2+x}$, determine the limit of $(u_n)$.
\item Prove that $0\le2-u_{n+1}\le\frac{2-u_n}{2}$ for every $n\in\N$. Deduce an upper bound for $2-u_n$ in terms of $n$, then an integer $N$ such that $0\le2-u_n\le10^{-3}$ for every $n\ge N$.\end{enumerate}

\ifmxteacher\mxnewpage\section{Teaching notes and solutions}\begin{teachernote}[Teacher notes]Test 2 comes before S19 and S20: assess in it the definitions and the operations already met, without requiring monotone convergence, Bernoulli or the squeeze theorem. Draw horizontal bands around the limit and have students look for an index from which all the terms stay inside them. Require convergence before solving the fixed-point equation. Show $(-1)^n$ to separate ``bounded'' from ``convergent''.\end{teachernote}

\subsection{Solutions to the activity}\begin{enumerate}
\item $u_0=5$, $u_9=\num{2,3}$ and $u_{99}=\num{2,03}$. Since $\frac{3}{n+1}>0$ at every index, $u_n>2$: no term equals $2$.
\item $|u_n-2|<\num{0,01}\Longleftrightarrow \frac{3}{n+1}<\num{0,01}\Longleftrightarrow n+1>300\Longleftrightarrow n\ge300$. The first suitable index is therefore $N=300$. At index $299$, the gap is exactly $\num{0,01}$ and the strict inequality is not satisfied.
\item For $\varepsilon>0$, take $N=\lceil\frac{3}{\varepsilon}\rceil$, the smallest integer greater than or equal to $\frac{3}{\varepsilon}$. If $n\ge N$, then $n+1>\frac{3}{\varepsilon}$, hence $|u_n-2|<\varepsilon$. This choice works, even when it is not minimal.
\item $|v_n-2|=\frac{1}{n+1}\to0$, hence $v_n\to2$. But $v_0=3$, $v_1=\frac{3}{2}$ and $v_2=\frac{7}{3}$: we observe a decrease then an increase. This convergent sequence is not monotonic.
\end{enumerate}

\paragraph{E1.}For $n\ge1$, $\frac{4n-1}{n+2}=\frac{4-\frac{1}{n}}{1+\frac{2}{n}}\to4$, since the denominator tends to $1\ne0$. Similarly $n^2-7n=n^2(1-\frac{7}{n})$; the first factor tends to $+\infty$ and the second to $1>0$, so the product tends to $+\infty$.

\paragraph{E2.}The conjugate gives $\sqrt{n+1}-\sqrt{n}=\frac{(n+1)-n}{\sqrt{n+1}+\sqrt n}=\frac{1}{\sqrt{n+1}+\sqrt n}$. The denominator tends to $+\infty$, so the difference tends to $0$.

\paragraph{E3.}For $n\ge1$, $-1\le\cos n\le1$ and $n>0$, so $-\frac{1}{n}\le u_n\le\frac{1}{n}$. Both bounds tend to $0$; the squeeze theorem gives $u_n\to0$.

\paragraph{E4.}The sequence is increasing and bounded above by $5$, so it converges. By passing to the limit in $u_n\le5$, its limit satisfies $\ell\le5$. It is not necessarily equal to $5$: $u_n=2-\frac{1}{n+1}$ is increasing, bounded above by $5$ and converges to $2$.

\paragraph{E5.}Since $|-\num{0,6}|<1$, $(-\num{0,6})^n\to0$, so $u_n\to3$. However $u_0=1$, $u_1=\num{4,2}$ and $u_2=\num{2,28}$: the sequence rises then falls and is not monotonic.

\subsection{Solutions to Homework 2}\begin{enumerate}
\item Let $P(n)$ be the statement ``$u_n$ is defined and $0\le u_n\le2$''. It is true at index $0$ since $u_0=0$. Fix $n\ge0$ and assume $P(n)$ is true. We have $2\le2+u_n\le4$, so $u_{n+1}=\sqrt{2+u_n}$ is defined and $\sqrt2\le u_{n+1}\le2$. Thus $P(n+1)$ is true. By induction, all the terms belong to $[0,2]$.
\item Let $R(n)$ be the statement ``$u_{n+1}\ge u_n$''. For $n=0$, $u_1=\sqrt2\ge0=u_0$. Fix $n\ge0$ and assume $R(n)$ is true. The fact that $x\mapsto\sqrt{2+x}$ is increasing on $[0,2]$ gives
\[u_{n+2}=\sqrt{2+u_{n+1}}\ge\sqrt{2+u_n}=u_{n+1}.\]
So $R(n+1)$ is true. By induction, $(u_n)$ is increasing. It is bounded above by $2$, so it converges to a real number $\ell\in[0,2]$.
\item The function $x\mapsto\sqrt{2+x}$ is continuous at $\ell$. By passing to the limit, $\ell=\sqrt{2+\ell}$. Since $\ell\ge0$, this is equivalent to $\ell^2-\ell-2=0$, that is $(\ell-2)(\ell+1)=0$. The root $-1$ is excluded; hence $\ell=2$.
\item Using the conjugate and $u_n\le2$,
\[0\le2-u_{n+1}=\frac{2-u_n}{2+\sqrt{2+u_n}}\le\frac{2-u_n}{2}.\]
Let $T(n)$ be the statement ``$0\le2-u_n\le\frac{2}{2^n}$''. At index $0$, $2-u_0=2$. If $T(n)$ is true for a fixed integer $n\ge0$, then
\[0\le2-u_{n+1}\le\frac12\frac2{2^n}=\frac2{2^{n+1}}.\]
The base case and the inductive step prove $T(n)$ at every index. Since $\frac{2}{2^{11}}=\frac{1}{1024}<10^{-3}$, $n=11$ suffices. This is an index guaranteed by an upper bound, with no claim of minimality. This question can be done before question 3: the number $2$ already appears in the bounds.
\end{enumerate}

\subsection{Solutions to the textbook or past-paper work}\paragraph{Exercise 17.} The quotients are at first indeterminate forms. For $n\ge1$,
\[u_n=\frac{2n-4}{n^2+1}=\frac{2\left(1-\frac2n\right)}{n\left(1+\frac1{n^2}\right)}\longrightarrow0.\]
Similarly,
\[v_n=\frac{2-5n}{4n+7}=\frac{\frac2n-5}{4+\frac7n}\longrightarrow-\frac54,\qquad w_n=n\frac{-1-\frac{10}{n^2}+\frac4{n^3}}{2+\frac3n+\frac1{n^2}}\longrightarrow-\infty.\]
The last quotient tends to $-\frac12$, and the factor $n$ to $+\infty$.
\paragraph{Exercise 24.} For every $n\ge1$, $-1\le(-1)^n\le1$. Since $n^2>0$, we obtain $-\frac1{n^2}\le u_n\le\frac1{n^2}$. Both bounds tend to $0$; the squeeze theorem gives $u_n\to0$.

\fi

